% ===== PRÉAMBULE CANONIQUE — à coller VERBATIM en tête de CHAQUE .tex =====
\documentclass[11pt,a4paper]{article}
\usepackage[utf8]{inputenc}\usepackage[T1]{fontenc}\usepackage{lmodern}
\usepackage[french]{babel}
\usepackage{amsmath,amssymb,mathtools}\usepackage{icomma}
\usepackage{siunitx}\sisetup{locale=FR}  % \qty{}{}, \si{}, \ang{}
\usepackage{xcolor}\usepackage{colortbl}
\usepackage{tikz}\usepackage{pgfplots}\pgfplotsset{compat=1.18}
\usepackage[siunitx,europeanresistors]{circuitikz} % circuits (élec.)
\usepackage[most]{tcolorbox}\usepackage{enumitem}\usepackage{array,booktabs}
\usepackage[a4paper,margin=2.2cm]{geometry}
\usetikzlibrary{arrows.meta,calc,angles,quotes,positioning,patterns,%
  backgrounds,shapes.geometric,decorations.pathmorphing,decorations.markings,intersections}
\definecolor{primary}{RGB}{222,49,99}\definecolor{primarydark}{RGB}{150,22,55}
\definecolor{defcol}{RGB}{0,114,178}\definecolor{methcol}{RGB}{52,92,148}
\definecolor{excol}{RGB}{0,135,90}\definecolor{attncol}{RGB}{200,40,55}
\tcbset{boxbase/.style={enhanced,breakable,boxrule=0.9pt,arc=2.5pt,
  left=3mm,right=3mm,top=2.3mm,bottom=2.3mm,fonttitle=\bfseries,coltitle=white}}
% --- boîtes TUTORIEL (noms EXACTS, ne pas renommer) ---
\newtcolorbox{cle}[1][]{boxbase,colback=defcol!6,colframe=defcol,
  title={\textsc{À retenir}\ifx\relax#1\relax\relax\else\textmd{ : \itshape #1}\fi}}
\newtcolorbox{methode}[1][]{boxbase,colback=methcol!7,colframe=methcol,sharp corners,
  title={\textsc{Méthode}\ifx\relax#1\relax\relax\else\textmd{ --- \itshape #1}\fi}}
\newtcolorbox{exemplebox}[1][]{boxbase,colback=excol!5,colframe=excol,
  title={\textsc{Exemple}\ifx\relax#1\relax\relax\else\textmd{ : \itshape #1}\fi}}
\newtcolorbox{piege}[1][]{boxbase,colback=attncol!6,colframe=attncol,
  title={\textsc{Piège}\ifx\relax#1\relax\relax\else\textmd{ : \itshape #1}\fi}}
\newtcolorbox{remarque}[1][]{boxbase,colback=black!4,colframe=black!45,
  title={\textsc{Remarque}\ifx\relax#1\relax\relax\else\textmd{ : \itshape #1}\fi}}
% --- boîtes EXERCICE / CORRIGÉ (noms EXACTS, auto-comptées) ---
\newtcolorbox[auto counter]{exo}[1][]{boxbase,colback=primary!4,colframe=primary,
  title={\textsc{Exercice}~\thetcbcounter\ifx\relax#1\relax\relax\else\textmd{ --- \itshape #1}\fi}}
\newtcolorbox[auto counter]{corrige}[1][]{boxbase,colback=excol!4,colframe=excol,
  title={\textsc{Corrigé}~\thetcbcounter\ifx\relax#1\relax\relax\else\textmd{ --- \itshape #1}\fi}}
\newcommand{\vv}[1]{\vec{#1}}\newcommand{\ds}{\displaystyle}
\newcommand{\R}{\mathbb{R}}\newcommand{\N}{\mathbb{N}}\newcommand{\Z}{\mathbb{Z}}
% =========================================================================
\begin{document}
\sisetup{per-mode=symbol}
% ================= CORRIGÉ FACILE =================

\begin{corrige}[calcul direct]
On applique directement la loi, en unités SI :
\[
F=\mathcal{G}\,\frac{m_A m_B}{d^2}
 =\num{6.67e-11}\times\frac{10\times 10}{(0{,}50)^2}
 =\num{6.67e-11}\times\frac{100}{0{,}25}
 =\num{6.67e-11}\times 400\approx\qty{2.7e-8}{\newton}.
\]
Une force minuscule : la gravitation entre objets courants est imperceptible.
\end{corrige}

\begin{corrige}[les deux forces]
Les deux forces forment un couple action--réaction.
\begin{enumerate}[label=\alph*)]
  \item Même intensité : $F_{B/A}=F_{A/B}=\qty{5.0e-8}{\newton}$.
  \item Elles sont de \textbf{sens opposés} (chaque corps attire l'autre vers lui).
  \item Elles sont portées par la droite $(AB)$ qui joint les centres des deux corps.
\end{enumerate}
\end{corrige}

\begin{corrige}[l'effet des masses]
$F$ est \emph{proportionnelle au produit} $m_A m_B$ (au numérateur).
\begin{enumerate}[label=\alph*)]
  \item $m_A\to 2m_A$ : la force double, $F=2F_0$.
  \item $m_B\to 3m_B$ : la force triple, $F=3F_0$.
  \item les deux doublent : $F=(2\times 2)F_0=4F_0$.
\end{enumerate}
\end{corrige}

\begin{corrige}[l'effet de la distance]
$F$ est \emph{inversement proportionnelle au carré} de $d$ (loi en $1/d^2$).
\begin{enumerate}[label=\alph*)]
  \item $d\to 2d$ : $F\to F_0/2^2=F_0/4$.
  \item $d\to 3d$ : $F\to F_0/3^2=F_0/9$.
  \item $d\to d/2$ : $F\to F_0/(1/2)^2=4F_0$.
\end{enumerate}
\end{corrige}

\begin{corrige}[force Terre--satellite]
\[
F=\mathcal{G}\,\frac{M_T\,m}{d^2}
 =\frac{\num{6.67e-11}\times\num{6.0e24}\times 500}{(\num{8.0e6})^2}
 =\frac{\num{2.0e17}}{\num{6.4e13}}\approx\qty{3.1e3}{\newton}.
\]
On a d'abord regroupé le numérateur $\num{6.67e-11}\times\num{6.0e24}=\num{4.0e14}$,
puis multiplié par $500$.
\end{corrige}

\end{document}
